% GENERATED FILE. Edit canonical lessons, then run npm run generate:books.
% canonical-source-hash: fnv1a64:fd499b6b9b56a673
% canonical-lessons: ZH-C94-dou, ZH-W94-dou, ZH-C94-yiyang, ZH-C94-nianqing, ZH-C94-pang, ZH-C94-shou, ZH-C94-practice

\chapter{All the Same}
\label{ch:zh-all-the-same}
\hlchaptermodality{94}

\begin{chapteropening}
\textbf{By the end of this chapter:} \emph{I can describe people and things, and say 'all of them'.}

Describing people, young, fat, thin, the same or not, and saying all of them with \zh{都}, built on one new written character: \zh{都}.
\end{chapteropening}

\section[\texorpdfstring{\zh{都} (dōu)}{都 (dōu)}]{\zh{都} (dōu) — all; both}

\label{lesson:ZH-C94-dou}



\begin{quote}
A new chapter, and it opens with the word that gathers people together.

\begin{itemize}
\raggedright
  \item \emph{Recall:} say \zh{太大了} — \emph{tài dà le}, too big
  \item \emph{Recall:} say \emph{xié tài zāng le}, the shoes are too dirty — \textbf{R1}, the chapter you have finished
  \item \emph{Recall:} say \zh{大家} — \emph{dàjiā}, everyone
\end{itemize}
\end{quote}

\subsection*{You'll want to know: \zh{都}}
\textbf{\zh{都}} — \emph{dōu} — \textquotedblleft{}all; both\textquotedblright{}.

\begin{quote}
\textbf{\zh{我们都是学生}} \emph{wǒmen dōu shì xuésheng} — We're all students.
\end{quote}

\begin{quote}
\textbf{\zh{他们都很好}} \emph{tāmen dōu hěn hǎo} — They're all fine.
\end{quote}

The next lesson puts the new character in your hand.

\begin{cousinweb}
\textbf{\zh{都}} (\emph{dōu}) is inherited from Middle Chinese \emph{tu} (Guangyun, in Baxter's transcription), glossed 'all together', and also 'a capital city', the sense said dū today. The regular change from \emph{tu} gives dū; why the word for 'all' is said dōu is not settled.
\end{cousinweb}

\begin{grammarlens}[title={What you've built}]
\zh{都} goes in front of the verb or the adjective and looks back at the people before it: \zh{我们都}, all of us. With two people it means both.
\end{grammarlens}

\subsection*{Your turn}
\begin{itemize}
\raggedright
  \item \emph{Say it:} \emph{dōu}
  \item \emph{Say it:} \emph{dōu}, once more
  \item \emph{Say it:} \zh{我们都是学生} — \emph{wǒmen dōu shì xuésheng} — We're all students.
  \item \emph{Say it:} \zh{他们都很好} — \emph{tāmen dōu hěn hǎo} — They're all fine.
  \item \emph{Recall:} say \emph{dàjiā}, then \emph{dōu} again
\end{itemize}

\begin{tcolorbox}[breakable,colback=teal!4,colframe=teal!35!black,title={Before you move on}]
What does \zh{都} mean? (\textquotedblleft{}All; both\textquotedblright{}.) Where does it go? (\textbf{In front of the verb, after the people}.)
\end{tcolorbox}
\section[\texorpdfstring{\zh{都} (dōu)}{都 (dōu)}]{\zh{都} — \zh{者}, then the ear on the right}

\label{lesson:ZH-W94-dou}



\begin{quote}
The word first, then the character inside it.

\begin{itemize}
\raggedright
  \item \emph{Recall:} say \emph{dōu} — \textbf{R1}, one lesson back, before it decays
\end{itemize}

\begin{quote}
\zh{日} \zh{那} \zh{在}
\end{quote}

\emph{Read it:} each one aloud, with nothing said first — \emph{rì}, \emph{nà}, \emph{zài}
\end{quote}

\subsection*{Script — \zh{者}, then the ear on the right}
\hlblockfigure{\detokenize{figures/ZH-W94-dou-filmstrip.pdf}}{How \zh{都} is written, stroke by stroke}

\begin{quote}
\zh{都}
\end{quote}

\textbf{Two pieces, left before right, and every stroke in them is one you have made.}

\textbf{\zh{者}, eight strokes}: first the foot of \zh{在}, a short bar, a vertical through it, then the long bar; then a long left-falling stroke from the upper right; then \zh{日} under them.

\textbf{The ear piece, two strokes}, on the right, as \zh{那} ends: across, back, round and hooking; then the long vertical.

8 + 2 = ten.

\subsection*{Your turn}
\begin{itemize}
\raggedright
  \item \emph{Write it:} \zh{都}, all ten, with the model in front of you
  \item \emph{Cover:} write it again from memory, and count the pieces, not the strokes
\end{itemize}

\begin{tcolorbox}[breakable,colback=teal!4,colframe=teal!35!black,title={Before you move on}]
How many strokes in \zh{都}? (\textbf{Ten}.) Which character ends the left piece? (\textbf{\zh{日}}.) Where have you written the ear piece? (\textbf{On the right of \zh{那}}.)

Source: \href{https://raw.githubusercontent.com/chanind/hanzi-writer-data/68d10a4b21150cae5e1ebbd223eed289cf32d90c/data/\%E9\%83\%BD.json}{Hanzi Writer Data \zh{都}, pinned stroke paths and medians}.
\end{tcolorbox}
\section[yíyàng]{yíyàng — the same}

\label{lesson:ZH-C94-yiyang}



\begin{quote}
Before the new one: read a line with two characters you have written in the last two chapters, then the tone of \zh{一}.

\begin{itemize}
\raggedright
  \item \emph{Read it:} \zh{他们也都很好}, aloud — \emph{tāmen yě dōu hěn hǎo}, they're all fine too
  \item \emph{Recall:} say \zh{一个} — \emph{yí ge}, with \zh{一} rising
\end{itemize}
\end{quote}

\subsection*{You'll want to know: yíyàng}
\emph{yíyàng} — \textquotedblleft{}the same; alike\textquotedblright{}.

\begin{quote}
\emph{wǒmen bù yíyàng} — We're not the same.
\end{quote}

\begin{quote}
\emph{wǒmen yíyàng dà} — We're the same age.
\end{quote}

\zh{一} before the fall of yàng rises: \emph{yí}. Its second character comes later; for now it is a word for the ear.

\begin{cousinweb}
\emph{yíyàng} is built from yī 'one' and yàng 'a pattern; a kind': one pattern. yī is inherited from Middle Chinese \emph{'jit} (Guangyun, in Baxter's transcription); yàng has no reading in the Guangyun, so no dictionary we use records an older form of it, and its history is not settled.
\end{cousinweb}

\begin{grammarlens}[title={What you've built}]
In \emph{tāmen yě dōu hěn hǎo}, \zh{也} comes before \zh{都}: too, and all of them. \emph{yíyàng} + an adjective says how two things match: \emph{yíyàng dà}, the same size or the same age.
\end{grammarlens}

\subsection*{Your turn}
\begin{itemize}
\raggedright
  \item \emph{Say it:} \emph{yíyàng}
  \item \emph{Say it:} \emph{yíyàng}, once more
  \item \emph{Say it:} \emph{wǒmen bù yíyàng} — We're not the same.
  \item \emph{Say it:} \emph{wǒmen yíyàng dà} — We're the same age.
  \item \emph{Recall:} say \emph{yí ge}, then \emph{yíyàng}, and hear the same rise
\end{itemize}

\begin{tcolorbox}[breakable,colback=teal!4,colframe=teal!35!black,title={Before you move on}]
What does \emph{yíyàng} mean? (\textquotedblleft{}The same\textquotedblright{}.) Which comes first, \zh{也} or \zh{都}? (\textbf{\zh{也}: \zh{也都}}.)
\end{tcolorbox}
\section[niánqīng]{niánqīng — young}

\label{lesson:ZH-C94-nianqing}



\begin{quote}
Before the new one: last lesson's word, then two words about age and people.

\begin{itemize}
\raggedright
  \item \emph{Recall:} say \emph{yíyàng}, the same — \textbf{R1}, one lesson back
  \item \emph{Recall:} say \emph{niánjì}, age — \textbf{R2}
  \item \emph{Recall:} say \emph{lǎoshī}, a teacher — \textbf{R4}, recognition at distance
\end{itemize}
\end{quote}

\subsection*{You'll want to know: niánqīng}
\emph{niánqīng} — \textquotedblleft{}young\textquotedblright{}.

\begin{quote}
\emph{wǒ de lǎoshī hěn niánqīng} — My teacher is very young.
\end{quote}

Its second character comes later; for now it is a word for the ear.

\begin{cousinweb}
\emph{niánqīng} is built from nián 'year', the \zh{年} you write, and qīng 'light', not heavy: light in years. Both are inherited from Middle Chinese, \emph{nen} and \emph{khjieng} (Guangyun, in Baxter's transcription), \emph{khjieng} glossed 'light in weight'.
\end{cousinweb}

\begin{grammarlens}[title={What you've built}]
\emph{niánqīng} is said of grown people; a child is \emph{xiǎo}. \emph{lǎo}, old, is its opposite.
\end{grammarlens}

\subsection*{Your turn}
\begin{itemize}
\raggedright
  \item \emph{Say it:} \emph{niánqīng}
  \item \emph{Say it:} \emph{niánqīng}, once more
  \item \emph{Say it:} \emph{wǒ de lǎoshī hěn niánqīng} — My teacher is very young.
  \item \emph{Recall:} say \emph{niánjì}, then \emph{niánqīng} again
\end{itemize}

\begin{tcolorbox}[breakable,colback=teal!4,colframe=teal!35!black,title={Before you move on}]
What does \emph{niánqīng} mean? (\textquotedblleft{}Young\textquotedblright{}.) Which written character does it open with? (\textbf{\zh{年}}.)
\end{tcolorbox}
\section[pàng]{pàng — fat; plump}

\label{lesson:ZH-C94-pang}



\begin{quote}
Before the new one: last lesson's word, then an animal and something to eat.

\begin{itemize}
\raggedright
  \item \emph{Recall:} say \emph{niánqīng}, young — \textbf{R1}, one lesson back
  \item \emph{Recall:} say \emph{māo}, a cat — \textbf{R2}
  \item \emph{Recall:} say \emph{dàngāo}, a cake — \textbf{R2}
\end{itemize}
\end{quote}

\subsection*{You'll want to know: pàng}
\emph{pàng} — \textquotedblleft{}fat; plump\textquotedblright{}.

\begin{quote}
\emph{māo tài pàng le} — The cat's too fat.
\end{quote}

Its character comes later; for now it is a word for the ear.

\begin{cousinweb}
The character for \emph{pàng} is in the Guangyun as \emph{phanH} (Guangyun, in Baxter's transcription), glossed 'half of a sacrificial animal'. That reading would give pàn, not pàng, and the sense is not 'fat', so the history of \emph{pàng} is not settled.
\end{cousinweb}

\begin{grammarlens}[title={What you've built}]
\emph{pàng} is a word for people and animals, and said of a person it is blunt: a cat can hear \emph{tài pàng le} without offence.
\end{grammarlens}

\subsection*{Your turn}
\begin{itemize}
\raggedright
  \item \emph{Say it:} \emph{pàng}
  \item \emph{Say it:} \emph{pàng}, once more
  \item \emph{Say it:} \emph{māo tài pàng le} — The cat's too fat.
  \item \emph{Say it:} \emph{māo yǒudiǎnr pàng} — The cat's a bit plump.
  \item \emph{Recall:} say \emph{dàngāo}, then \emph{pàng} again
\end{itemize}

\begin{tcolorbox}[breakable,colback=teal!4,colframe=teal!35!black,title={Before you move on}]
What does \emph{pàng} mean? (\textquotedblleft{}Fat; plump\textquotedblright{}.) Say the cat is too fat. (\emph{māo tài pàng le}.)
\end{tcolorbox}
\section[shòu]{shòu — thin; slim}

\label{lesson:ZH-C94-shou}



\begin{quote}
Before the new one: last lesson's word, then tall and short.

\begin{itemize}
\raggedright
  \item \emph{Recall:} say \emph{pàng}, fat — \textbf{R1}, one lesson back
  \item \emph{Recall:} say \emph{gāo}, tall — \textbf{R3}, durable
  \item \emph{Recall:} say \emph{ǎi}, short — \textbf{R3}
\end{itemize}
\end{quote}

\subsection*{You'll want to know: shòu}
\emph{shòu} — \textquotedblleft{}thin; slim\textquotedblright{}.

\begin{quote}
\emph{tā hěn gāo, hěn shòu} — He's tall and thin.
\end{quote}

Its character comes later; for now it is a word for the ear.

\begin{cousinweb}
\emph{shòu} is inherited from Middle Chinese \emph{srjuwH} (Guangyun, in Baxter's transcription), glossed, through an older dictionary, 'lean'.
\end{cousinweb}

\begin{grammarlens}[title={What you've built}]
\emph{pàng} and \emph{shòu}, \emph{gāo} and \emph{ǎi}: two pairs that describe a body, each with a plain \zh{很}.
\end{grammarlens}

\subsection*{Your turn}
\begin{itemize}
\raggedright
  \item \emph{Say it:} \emph{shòu}
  \item \emph{Say it:} \emph{shòu}, once more
  \item \emph{Say it:} \emph{tā hěn gāo, hěn shòu} — He's tall and thin.
  \item \emph{Recall:} say \emph{pàng}, then \emph{shòu}, the pair
\end{itemize}

\begin{tcolorbox}[breakable,colback=teal!4,colframe=teal!35!black,title={Before you move on}]
What does \emph{shòu} mean? (\textquotedblleft{}Thin; slim\textquotedblright{}.) And its opposite? (\emph{\textbf{pàng}}.)
\end{tcolorbox}
\section[Practice]{\zh{我们都是学生}}

\label{lesson:ZH-C94-practice}



\begin{quote}
The chapter closes here, with a photograph of a class.

\begin{itemize}
\raggedright
  \item \emph{Recall:} say \emph{shòu}, thin — \textbf{R1}, one lesson back
  \item \emph{Recall:} say \zh{同学} — \emph{tóngxué}, a classmate — \textbf{R4}, recognition at distance
  \item \emph{Recall:} say \emph{gānjìng}, clean — \textbf{R2}
\end{itemize}
\end{quote}

\subsection*{The exchange}
\begin{quote}
\textbf{A:} \zh{他们都是你的同学吗}
\end{quote}

\begin{quote}
\emph{tāmen dōu shì nǐ de tóngxué ma} — Are they all your classmates?
\end{quote}

\begin{quote}
\textbf{B:} \zh{对} \zh{我们都是学生}
\end{quote}

\begin{quote}
\emph{duì, wǒmen dōu shì xuésheng} — Yes, we're all students.
\end{quote}

\begin{quote}
\textbf{A:} \emph{nǐmen dōu hěn niánqīng} — You're all so young.
\end{quote}

\begin{quote}
\textbf{B:} \emph{wǒmen yíyàng dà} — We're the same age.
\end{quote}

\begin{quote}
\textbf{A:} \zh{那个男人呢}
\end{quote}

\begin{quote}
\emph{nàge nánrén ne} — And that man?
\end{quote}

\begin{quote}
\textbf{B:} \emph{tā shì wǒmen de lǎoshī, tā hěn gāo, hěn shòu} — He's our teacher. He's tall and thin.
\end{quote}

\begin{quote}
\textbf{A:} \emph{nǐ de māo hěn pàng} — Your cat's fat.
\end{quote}

\begin{quote}
\textbf{B:} \zh{也很大}
\end{quote}

\begin{quote}
\emph{yě hěn dà} — And big.
\end{quote}

\zh{都} gathers everybody before it into one: \zh{他们都}, \zh{我们都}, \zh{你们都}. \emph{yíyàng} matches them up; \emph{gāo}, \emph{shòu} and \emph{pàng} tell them apart.

\begin{tcolorbox}[breakable,skin=enhanced,colback=orange!6,colframe=orange!45!black,boxrule=0.5pt,arc=1mm,left=6pt,right=6pt,top=4pt,bottom=4pt,fonttitle=\bfseries,title={Read this}]
\emph{Read it:} the exchange through once without stopping, then answer the question below
\end{tcolorbox}

\begin{grammarlens}[title={What you've built — describing people}]
\begin{itemize}
\raggedright
  \item \textbf{\zh{我们都是学生}} \emph{wǒmen dōu shì xuésheng} — we're all students.
  \item \emph{wǒmen yíyàng dà} — we're the same age.
  \item \emph{tā hěn gāo, hěn shòu} — he's tall and thin.
  \item \textbf{\zh{他们也都很好}} \emph{tāmen yě dōu hěn hǎo} — they're all fine too.
\end{itemize}

\textbf{Listen.} \emph{Nǐmen dōu shì xuésheng ma?}

\textbf{Speak.} Describe two people you know, then say what they have in common with \zh{都}.

\emph{Write it:} \zh{都} — the left piece first, closed by its \zh{日}, then the ear on the right
\end{grammarlens}

\subsection*{Your turn}
\begin{itemize}
\raggedright
  \item \emph{Say it:} \zh{我们都是学生} — \emph{wǒmen dōu shì xuésheng} — We're all students.
  \item \emph{Say it:} \zh{他们也都很好} — \emph{tāmen yě dōu hěn hǎo} — They're all fine too.
  \item \emph{Say it:} the whole exchange, both parts
  \item \emph{Write it:} \zh{我们都是学生}
\end{itemize}

\begin{tcolorbox}[breakable,colback=teal!4,colframe=teal!35!black,title={Before you move on}]
Say you are all students. (\textbf{\zh{我们都是学生}}.) Say you are the same age. (\emph{wǒmen yíyàng dà}.) Say somebody is young. (\emph{tā hěn niánqīng}.) Which comes first, \zh{也} or \zh{都}? (\textbf{\zh{也}}.)
\end{tcolorbox}
